\section{Skill Learning and Self-Healing}
\label{sec:skills}

The skill system is the core differentiator of OpenSkynet. Skills are named, versioned, executable recipes consisting of textual steps, structured browser actions (navigate, click, type, scroll, extract), and rich metadata (verification criteria, retry policy, timeout, allowed tools, input/output signatures). The full lifecycle is shown in Figure~\ref{fig:skill_lifecycle}.

\begin{figure}[t]
    \centering
    \begin{tikzpicture}[
        x=2.4cm, y=1.8cm,
        every node/.style={font=\sffamily, rounded corners=4pt},
        st/.style={
            minimum height=1.7cm, minimum width=2.0cm,
            text=white, font=\sffamily\scriptsize, align=center,
            fill, draw=none,
        },
        arr/.style={-{Stealth[length=4pt, width=3pt]}, color=cArrow, very thick},
        lbl/.style={font=\sffamily\tiny, text=cBorder, fill=white, inner sep=1.5pt},
        fb/.style={-{Stealth[length=4pt, width=3pt]}, color=cArrow, very thick, dashed},
    ]
        \node[st, fill=cSkills]   at (0,0) (s1) {\faVideo\\\textbf{1. Record}\\[2pt]{\tiny}Capture trajectory\\(actions + DOM + screenshots)};
        \node[st, fill=cSubagent] at (2.8,0) (s2) {\faMagic\\\textbf{2. Extract}\\[2pt]{\tiny}LLM review +\\3-cycle refinement};
        \node[st, fill=cManager]  at (5.6,0) (s3) {\faPlayCircle\\\textbf{3. Execute}\\[2pt]{\tiny}Programmatic or\\LLM-driven};
        \node[st, fill=cHealing]   at (8.4,0) (s4) {\faMedkit\\\textbf{4. Self-Heal}\\[2pt]{\tiny}Screenshot + DOM\\$\to$ auto-patch};
        \node[st, fill=cMemory]   at (11.2,0) (s5) {\faClipboardCheck\\\textbf{5. Audit}\\[2pt]{\tiny}Staleness detect\\$\to$ archive / prune};

        \draw[arr] (s1.east) -- (s2.west);
        \draw[arr] (s2.east) -- (s3.west);
        \draw[arr] (s3.east) -- (s4.west);
        \draw[arr] (s4.east) -- (s5.west);

        \draw[fb, color=cHealing]
            (s4.north) -- ++(0, 0.5) -|
            node[lbl, above, pos=0.25] {\faRedo~~retry with healed steps}
            (s3.north);

        \draw[fb, color=cMemory]
            (s5.south) -- ++(0, -0.5) -|
            node[lbl, below, pos=0.25] {\faSync~~refresh stale skills}
            (s1.south);
    \end{tikzpicture}
    \caption{Skill lifecycle: trajectories are recorded, skills are extracted and refined, execution heals on failure, and periodic auditing prunes stale skills.}
    \label{fig:skill_lifecycle}
\end{figure}

\subsection{Skill Storage and Versioning}

Skills are stored as dual-format files on disk: \texttt{skill.json} (structured, machine-readable) and \texttt{SKILL.md} (human-readable Markdown with YAML frontmatter). The Skill Engine supports two storage scopes: a global directory (\texttt{\textasciitilde/.openskynet/skills/}) and a project-local directory (\texttt{.openskynet/skills/}), with project-level skills taking priority. Every update creates a versioned archive in a \texttt{history/} subdirectory, enabling rollback to any prior version. A lock file (\texttt{skills-lock.json}) tracks installation provenance, version hashes, and update detection for skills installed from the centralized skills hub or GitHub repositories.

\subsection{Skill Learning Pipeline}

Skill learning occurs automatically after each successful browser task. The \texttt{SkillLearnerAgent} orchestrates a multi-stage pipeline:

\begin{enumerate}[leftmargin=*,itemsep=1pt]
    \item \textbf{Pre-check:} If the task was successful and $\geq$2 browser actions were performed, check for similar existing skills using vector similarity search. If a close match exists (cosine similarity $\geq$ 0.1), skip LLM review and directly decide to patch the existing skill.
    \item \textbf{LLM Review:} A structured prompt (loaded from a configurable template) asks the LLM to evaluate three criteria: (a) count of actions $\geq$ 5, (b) task reusability, and (c) novelty relative to existing skills. At least 2 of 3 must be satisfied to proceed.
    \item \textbf{Iterative Refinement:} The initial skill draft is refined through up to 3 cycles of LLM critique. Each cycle feeds the current draft back to the model with instructions to critique and improve step specificity, add CSS selectors, identify pitfalls, and suggest verification criteria.
    \item \textbf{Failure-Aware Context:} Recent failed trajectories (up to 5) are loaded from the trajectory database and included as contextual examples. The LLM is instructed to avoid patterns that led to prior failures.
    \item \textbf{Similar-Skill Merging:} If a similar skill exists and \texttt{should\_patch} is set, the new steps are merged into the existing skill rather than creating a duplicate. If no similar skill exists, a new skill is created.
    \item \textbf{Security Gate:} Before saving, all skill content (name, description, steps) is scanned for prompt injection patterns, credential leaks, and destructive commands.
\end{enumerate}

Skills can also be learned from screen recordings via the \texttt{TraceToSkill} module. This module extracts up to 25 key frames from a recorded session, overlays cursor positions on each frame, and sends the multimodal sequence (frames + action log + DOM summaries) to an LLM. A two-phase post-processing step extracts reusable \texttt{\{\{variable\}\}} placeholders and refines steps with DOM-specific selectors for robustness.

\subsection{Self-Healing Execution}

When a skill execution fails---whether due to a changed DOM structure, a missing element, an unexpected popup, or a verification failure---the Healer module is invoked automatically:

\begin{enumerate}[leftmargin=*,itemsep=1pt]
    \item \textbf{Evidence Capture:} A screenshot of the browser state at failure and a truncated DOM snapshot (up to 3000 characters) are captured. The screenshot is limited to 5MB to prevent base64 overflow.
    \item \textbf{Multimodal Diagnosis:} The error context, original skill steps, screenshot (as base64-encoded image), and DOM snapshot are sent to an LLM with a structured healer prompt. The model returns a JSON object containing the identified failing step, root cause, corrected steps, and a confidence level (low/medium/high).
    \item \textbf{Unfixable Detection:} If the LLM responds with an explicit \texttt{error} key, the system marks the skill as unfixable and terminates without patching, preventing cascading failures.
    \item \textbf{Auto-Patching:} The corrected steps are applied to the skill via the engine's \texttt{patch()} method, which creates a new version and archives the previous one. The confidence and reasoning are logged for audit.
    \item \textbf{Re-execution:} The healed skill is re-executed. If healing succeeds, the skill's success rate metric is updated. If it fails again, the skill is marked for manual review.
\end{enumerate}

This creates a \emph{continuous improvement loop}: skills become more robust over time through actual usage, without requiring human intervention to update selectors or steps when web pages change.

\subsection{Skill Audit}

A periodic \texttt{SkillAuditor} reviews installed skills for staleness. Skills unused for 30+ days are flagged; the auditor consults an LLM for additional nuance (e.g., detecting truly broken vs.\ temporarily unused skills) and recommends actions: archive (set category to ``archived'') or delete. The auditor also identifies skills with low user preference scores (from the DPO-inspired preference learner) and suggests pruning.
